\scanpage{201}
\fbox{Sous-groupe $G^\tau$}

Soit $G$ un schéma en groupes sur le \uline{corps} $k$,
localement de type fini sur $k$. Alors :

\deleted{1°) La composante connexe de l'élément neutre $e$
est de type fini sur $k$, et \unsure{relt.|ablt.} irréductible}

\deleted{\unsure{2}}1°) Les composantes connexes \inserted{$G_i$} de $G$ sont de type
fini sur $k$, et irréductibles

2°) La composante connexe \inserted{$G^{\deleted{\unsure{0}}}_0$} de $e$ est
\unsure{relt.|ablt.|géom.} irréductible, et c'est un sous-groupe invariant

3°) \deleted{Si $k$ est alg. clos, alors les \unsure{$G_i$} \unsure{ouvert et fermé} les composantes
connexes de $G$ sont les parties de la forme
$aG^{\unsure{0|\circ}}_0$, où \unsure{$a$} $a\in G(k)$} \inserted{Les $G_i$ sont stables sous $G_0$. Si $G_i$ contient un pt $a$ rat. sur $k$, alors $G_i=G_0a$}

4°) Le \deleted{groupe}\inserted{schéma quotient} $G/G^{\deleted{\unsure{0|\circ}}}_0$ existe il est discret i.e.
localement fini sur $k$, $G\unsure{-|\to}G/G^{\deleted{\unsure{0|\circ}}}_0$ est
un morphisme \deleted{\unsure{plat}} fid plat. $G/G_0$ est
un $k$ groupe, sépa\uuline{rable} et loc. fini. \deleted{\unsure{sur $k$.}}

\deleted{La composante connexe $G^0$ de $e_{\bar k}$ dans $G_{\bar k}$ est}
Supposons d'abord $k$ alg. clos. \deleted{Soit $G_0$}
\deleted{définie sur} Comme $G_0$ est connexe,
il en est de $=$ de $G_0\times G_0$ (\uuline{référence} \ldots)
% slanted
\margin{$\times$ Attention, $G_0$ n'est pas à \uline{priori} de type fini.}
donc $G_0\times G_{\unsure{0|e}}\unsure{-|\to}G$ induit $G_0\times G_0-G_0$,
i.e. $G_0$ est un sous-groupe. Soit $U$
un voisinage ouvert \inserted{\unsure{affine}} de $e$, je dis que
$U^2=G_0$ \deleted{\unsure{et}} donc $G_0$ est \unsure{quasi|qu.} cpct et
% archive sheet number: 200
